\section{Conclusion}
\label{sec:conclusion}

We introduced weakly non-separable Hamiltonians and a fixed-point reduction
that converts them into a sequence of separable, GAN-compatible MFG
subproblems. The central contribution is an exact outer-map $L^2$
contraction theorem (\Cref{thm:main}) yielding $\mathcal{W}_2$ convergence
(\Cref{cor:W2}) under an explicit weak-coupling condition. In the
density-only subcase, the contraction factor reduces to the clean rate
$\eps L_R/\lambda$, and this threshold is sharp for the linearized ergodic
model on the torus. For residuals that also affect the velocity field, the
rate contains the additional constants $L_{R,p}$ and $L_{p\rho}$, which
account for momentum variation and frozen-density dependence of
$\nabla_pR$.

\paragraph{Limitations.}
The contraction theorem covers densities bounded away from vacuum
($\rho \geq \rho_{\min} > 0$), and the positivity condition
$c_0\rho_{\min}-3\eps L_{R,p}\rho_{\max}>0$ depends on the density contrast;
problems with vacuum regions or very sharp density gradients are outside
the present analysis. The numerical validation is conducted in the
linearized regime around the uniform equilibrium, where the density-only
and traffic linearizations allow the predicted rates to be measured
directly. These experiments test the outer fixed-point mechanism rather
than a full stochastic adversarial training pipeline.
The analysis also assumes exact inner solves; the bias from finite
inner iterations is bounded in \Cref{prop:inexact}, but a tighter treatment
of optimization and sampling error remains open.

\paragraph{Future directions.}
Several extensions are natural.
The most important next step is a full neural implementation in which
APAC-Net or another adversarial MFG solver is used as the inner routine,
followed by pipeline-level comparisons with MFDGM, PINN/DGM solvers, and
finite-difference references in low dimension.
When the weak-coupling condition fails, hybrid schemes that switch from
distributional inner solvers to residual-minimization methods such as MFDGM
may provide a practical route to more strongly non-separable models.
Extension to mean field control problems \citep{carmona2018probabilistic},
where the social planner optimizes over the whole population, is also
of interest.
Finally, applying the fixed-point reduction to large-scale multi-agent
reinforcement learning and to engineering systems modeled by
congestion-type MFGs are promising practical directions.

\paragraph{Code and data availability.}
The full LaTeX source, the Python code for all experiments, and the
generated figures are available at
\url{https://github.com/hikmatulloismatov1999/MFG-GAN-Paper}.
Each experiment runs in under one minute on a standard laptop using only
NumPy and Matplotlib; no GPU or specialised hardware is required to
reproduce the tables and figures in \Cref{sec:experiments}.

\paragraph{Acknowledgements.}
The authors thank the Mean-Field Games and Numerical Analysis groups at
the King Abdullah University of Science and Technology (KAUST) for
helpful discussions.
